\scanpage{29}
\inserted{\deleted{Moyennant conjecture de Hodge,}} On \unsure{récupère} $\mathcal M$ comme $\mathrm{Rep}(G)$\inserted{(cf a))}, $T_B$ \inserted{(sans structure de Hodge)} comme le foncteur oubli sur $\mathrm{Rep}(G)$, $T_{\mathrm{DR}}$ comme $\deleted{T_B}P\overset{G}{\deleted{\unsure{\otimes|\circ}}\times}(T_D\otimes_{\mathbb Q}k)$ (cf e), f)), \deleted{la structure} de sorte que $G'\simeq\underline{\mathrm{Aut}}_{\unsure{\ominus|\otimes}}(T_{\mathrm{DR}})$.

\deleted{Alors $Q$ est un bitorseur \unsure{sous|sur} $(G',G'')$ (cf b)c)d))}

Alors on a $H'\subset G'$, un bitorseur $Q$ sur $(H',H'')$, \deleted{deux \unsure{hom.} $i_1,i_2$ \unsure{d|et} \unsure{avec}} et $T_{\mathrm{Hdg}}\simeq T_{\mathrm{DR}}\overset{H'}{\times}Q$, enfin \deleted{$i_1,i_2\colon G_m\mathrel{-}$} $G''\simeq\underline{\mathrm{Aut}}_{\otimes}(T_{\mathrm{Hdg}})$ et $i_1,i_2\colon \unsure{\mathbb G_{m_k}|\mathbb G_{m\,k}}\mathrel{-}G''$ définissant une bigraduation sur $T_{\mathrm{Hdg}}$. Les relations \deleted{\unsure{propr.|prop.}} définissant $H''$ en termes de $i_1$ impliquent que cette bigraduation correspond à une \deleted{\unsure{graduation}} \uline{filtration} de $T_{\mathrm{DR}}$. On en conclut une filtration sur $T_{\mathrm{DR}}\otimes_k\mathbb C$. Or $\varphi$ définit \unsure{un|une} isom. de ce dernier avec $T_B\otimes_{\mathbb Q}\mathbb C$, qui se trouve donc muni d'une filtration.

\uline{Conditions} à imposer aux données a) b) c) d) e) f) g)

\deleted{\illegible{1}} 1) Pour tte repr. (ou une repr. fidèle) $V$ de $G$, la filtration qu'on obtient ainsi sur $V_{\mathbb C}$ \deleted{satisfait} fait de $V$ un module de Hodge\inserted{(condition ouverte)}. 2) Ce dernier est \uuline{motivique}

3) La restriction des scalaires de $\mathbb C$ à $k$ de $V_{\mathbb C}$, déduite de $\varphi\colon T_{\mathrm{DR}}\otimes_k\mathbb C\simeq T_D\otimes_{\mathbb Q}\mathbb C$, correspond à une restriction des scalaires du motif sur $\mathbb C$ défini par la structure de Hodge, grâce à 1). \uline{NB} On ne sait pas expliciter \unsure{2)|3)} et 3) \deleted{\unsure{en termes}} directement en termes des données a) .... !

On trouve alors, \inserted{\uline{si on admet les conjectures de Hodge,}} un hom
% paper watermark, sideways
{\centering \textsc{\unsure{Extra Strong}}\\
\textit{\unsure{Ma}\unsure{in|nn}\unsure{ol|on}\unsure{uy|ny|y} \unsure{1808}}\par}
% archive sheet number: 28
